\documentclass[11pt,a4paper]{article}

\usepackage[margin=23mm]{geometry}
\usepackage{amsmath,amssymb,mathtools,bm}
\usepackage{booktabs,longtable,array}
\usepackage{algorithm}
\usepackage{algpseudocode}
\usepackage{xurl}
\usepackage[hidelinks]{hyperref}
\usepackage{enumitem}
\usepackage{microtype}

\newcommand{\Ftwo}{\mathbb F_2}
\newcommand{\supp}{\operatorname{supp}}
\newcommand{\TopK}{\operatorname{TopK}}
\newcommand{\N}{\mathcal N}
\newcommand{\Bcal}{\mathcal B}
\newcommand{\Crel}{\mathcal C_{\mathrm{rel}}}
\newcommand{\Vrel}{\mathcal V_{\mathrm{rel}}}
\newcommand{\Vphys}{\mathcal V_{\mathrm{phys}}}
\newcommand{\Cphys}{\mathcal C_{\mathrm{phys}}}
\newcommand{\xor}{\mathbin{\oplus}}
\newcommand{\sgn}{\operatorname{sgn}}
\newcommand{\clip}{\operatorname{clip}}
\newcommand{\ind}{\mathbf 1}
\newcommand{\argmax}{\operatorname*{arg\,max}}

\title{Adaptive Failure-Localized Graph Refactorization for BP Decoding\\
\large Algorithm Specification and Initial Implementation Design}
\author{}
\date{Draft specification, September 2026}

\begin{document}
\maketitle

\begin{abstract}
This note specifies an initial adaptive graph-refactorization decoder for binary
syndrome-decoding problems. The decoder alternates between belief-propagation
(BP) decoding and exact, local factor-graph transformations. After a failed BP
run, a set of suspicious physical variables is selected from BP-derived failure
weights. Biclique candidates are then enumerated from pairs of relevant check
nodes. A selected biclique is factorized by introducing one auxiliary variable
and one auxiliary check, preserving the physical solution set exactly while
changing the factorization seen by BP. Two factorization-selection policies
are defined: an adaptive failure-weighted net 4-cycle score and a Shen-style
greedy internal 4-cycle-count baseline. Multiple transformations are performed
sequentially with candidate recomputation after every accepted transformation,
which removes ambiguity caused by overlapping bicliques. Three BP variants are
specified: synchronous parallel Min-Sum, horizontal serial Min-Sum, and the
qDither algorithm of Cai et al. The purpose of this document is to define the
first implementation unambiguously; it does not claim that the proposed
adaptive score is optimal.
\end{abstract}

\tableofcontents
\newpage

\section{Scope and design principles}

The decoder considered here does not search directly over complete corrections.
Its adaptive degree of freedom is the \emph{factorization} used by BP. Every
graph transformation must preserve the original physical decoding problem
exactly. The graph may acquire auxiliary variables and auxiliary checks, but
the set of valid physical corrections must remain unchanged.

The first implementation is deliberately restricted to biclique
factorizations. Arbitrary linear refactorizations, explicit optimization over
larger check subsets, 6-cycle-specific transformations, and joint optimization
of several transformations are outside the present specification.

The complete decoder has the following structure:
\[
\boxed{
\text{BP}
\;\longrightarrow\;
\text{failure weighting}
\;\longrightarrow\;
\text{biclique discovery}
\;\longrightarrow\;
\text{bounded exact factorization}
\;\longrightarrow\;
\text{BP on the new graph}.
}
\]
The BP--factorization relay is repeated a bounded number of times.

\section{Original decoding problem and current factorization}

\subsection{Physical decoding problem}

Let
\[
H_0\in\Ftwo^{m_0\times n_0},\qquad
s_0\in\Ftwo^{m_0}
\]
denote the immutable physical decoding matrix and measured syndrome. The
physical fault vector is
\[
x\in\Ftwo^{n_0},
\]
and a valid physical correction satisfies
\[
H_0x=s_0.
\]

For physical variable $j$, let $p_j$ be its prior fault probability and
define the base LLR
\[
\lambda_j^{\mathrm{base}}
=
\log\frac{1-p_j}{p_j}.
\]
Positive LLR favors the value $0$.

The original physical Tanner graph is denoted
\[
G_0=(\Vphys,\Cphys,E_0).
\]

\subsection{Current augmented factor graph}

At graph-relay stage $r$, BP runs on
\[
G^{(r)}=(V^{(r)},C^{(r)},E^{(r)})
\]
with binary linear constraints
\[
H^{(r)}z=s^{(r)}.
\]
The current variable vector $z$ contains all physical variables and possibly
auxiliary variables.

Every current variable $v\in V^{(r)}$ is assigned a physical-support vector
\[
P(v)\in\Ftwo^{n_0}
\]
such that the value of $v$ is the parity
\[
z_v=P(v)^Tx.
\]
For an original physical variable $x_j$,
\[
P(x_j)=e_j,
\]
where $e_j$ is the $j$-th standard basis vector.

For an auxiliary variable $y$ defined as
\[
y=\bigoplus_{v\in S} z_v,
\]
its physical support is
\[
P(y)=\bigoplus_{v\in S}P(v).
\]
The XOR here is essential: if a physical variable occurs an even number of
times through nested auxiliary definitions, it cancels.

The central invariant of every graph transformation is
\[
\boxed{
\{x:H_0x=s_0\}
=
\left\{
x:\exists!\,z_{\mathrm{aux}}\;
H^{(r)}z=s^{(r)}
\right\}.
}
\]
Thus each valid physical assignment has exactly one compatible assignment of
the current auxiliary variables.

\subsection{Unary priors}

The immutable unary prior for a physical variable remains
\[
\lambda_v^{\mathrm{base}}=\lambda_j^{\mathrm{base}}
\qquad (v=x_j).
\]
An auxiliary variable represents a deterministic parity and receives no
independent physical evidence:
\[
\boxed{
\lambda_y^{\mathrm{base}}=0.
}
\]
This zero LLR is a neutral unary factor; it must not be replaced by an
inherited BP marginal as a permanent prior. Inherited marginals are used only
to initialize message passing, as specified in Section~\ref{sec:handoff}.

\section{BP engine interface}

The graph-refactorization logic treats BP through a common interface.

\subsection{BP input and output}

A call
\[
\mathrm{RunBP}(G,\lambda^{\mathrm{base}},q^{\mathrm{init}},\Theta_{\mathrm{BP}})
\]
takes a current factor graph, its immutable unary LLR vector, an initialization
vector used only for initial variable-to-check messages, and BP-variant
parameters. The call returns
\[
\mathsf{BPResult}
=
(\mathsf{converged},\hat z,L^{\mathrm{final}},
\mathcal H_L,\mathsf{iterations}),
\]
where $L^{\mathrm{final}}$ is the final marginal LLR vector and $\mathcal H_L$
is the trailing marginal history required by the failure diagnostic.

Any candidate solution is accepted only after projecting $\hat z$ to its
physical coordinates $\hat x$ and checking
\[
H_0\hat x=s_0.
\]

\subsection{Common Min-Sum rule}

For the parallel and serial variants, and for the inner Min-Sum recursion of
qDither, the base check update is
\[
m_{a\to v}
=
\alpha_{\mathrm{MS}}
(-1)^{s_a}
\left(
\prod_{u\in N(a)\setminus\{v\}}
\sgn m_{u\to a}
\right)
\min_{u\in N(a)\setminus\{v\}}
|m_{u\to a}|,
\]
and the variable update is
\[
m_{v\to a}
=
\lambda_v^{\mathrm{base}}
+
\sum_{b\in N(v)\setminus\{a\}}
m_{b\to v}.
\]
The marginal is
\[
L_v
=
\lambda_v^{\mathrm{base}}
+
\sum_{b\in N(v)}m_{b\to v}.
\]
The normalization parameter satisfies $0<\alpha_{\mathrm{MS}}\le1$; setting
$\alpha_{\mathrm{MS}}=1$ gives ordinary Min-Sum.

\subsection{Variant P: parallel Min-Sum}

The parallel variant is synchronous. At iteration $t$, all check-to-variable
messages are computed from variable-to-check messages from iteration $t-1$;
then all variable-to-check messages are updated; then marginals and the hard
decision are computed. No message produced by one check during a given
check-update half-iteration is used by another check in that same
half-iteration.

\subsection{Variant S: horizontal serial Min-Sum}

The serial variant processes check nodes one at a time. Let
\[
\pi_t=(a_1,\ldots,a_{|C|})
\]
be the check order in iteration $t$. When check $a_k$ is processed, its new
check-to-variable messages are computed using the currently available
variable-to-check messages. The affected variable marginals and outgoing
variable-to-check messages are updated immediately, so later checks in the
same iteration can use the new information.

The serial-order option is
\[
\texttt{serial\_order}\in
\{\texttt{natural},\texttt{random\_per\_iteration}\}.
\]
The randomized option draws an independent permutation $\pi_t$ at the start of
each iteration using the decoder RNG seed. GARI uses a randomized horizontal
serial schedule in the serial portion of its hybrid schedule~\cite{gari}.

\subsection{Variant D: qDither Min-Sum}

The qDither variant follows Cai et al.~\cite{qdither}. A faithful paper-mode
run has two phases.

\paragraph{Phase I.}
Run standard syndrome-based Min-Sum from the base priors. If a valid solution
is found, return it.

\paragraph{Phase II.}
For chain $r=1,\ldots,n_{\mathrm{chain}}$, form a new bias
\[
\lambda_j^{(r)}
=
(1-\chi_j^{(r)})\lambda_j^{\mathrm{base}}
+
\chi_j^{(r)}q_j^{(r-1)},
\qquad
\chi_j^{(r)}\sim\mathrm{Unif}[\alpha_D,\beta_D].
\]
For each variable $j$, draw once per chain
\[
I_j^{(r)}\sim\mathrm{Bernoulli}(\rho_D).
\]
After the ordinary variable-to-check update produces a provisional
$m_{j\to a}^{(\ell)}$, define
\[
\delta_{j\to a}^{(\ell)}
=
\frac{
1-
\sgn m_{j\to a}^{(\ell)}
\sgn m_{j\to a}^{(\ell-1)}
}{2}.
\]
If $I_j^{(r)}=1$, replace the message by
\[
m_{j\to a}^{(\ell)}
\leftarrow
m_{j\to a}^{(\ell)}
+
\delta_{j\to a}^{(\ell)}
m_{j\to a}^{(\ell-1)}.
\]
After an unsuccessful chain, set
\[
q_j^{(r)}=L_j^{\mathrm{final}}
\]
and continue to the next chain. The qDither specification in Cai et al. uses
$q^{(0)}=\lambda^{\mathrm{base}}$.

The implementation option
\[
\texttt{qdither\_handoff}\in
\{\texttt{paper},\texttt{graph\_warm}\}
\]
is defined as follows. In \texttt{paper} mode, the paper initialization is used
exactly and cross-graph marginal inheritance is not injected into qDither
Phase I. In \texttt{graph\_warm} mode, the first qDither message state is
initialized from the graph-relay handoff vector defined in
Section~\ref{sec:handoff}; this is a graph-relay/qDither hybrid and must not be
reported as an exact reproduction of qDither.

\section{Failure weights and suspicious physical variables}

\subsection{Residual syndrome}

If BP fails, project the final hard decision to physical variables $\hat x$
and compute the residual in the immutable original basis:
\[
r^{\mathrm{phys}}
=
s_0\oplus H_0\hat x.
\]
Let
\[
\mathcal R
=
\{a\in\Cphys:r_a^{\mathrm{phys}}=1\}.
\]

\subsection{Uncertainty and oscillation}

For physical variable $j$, let the last $W$ physical marginal LLRs of the
failed BP call be $L_j^{(t)}$. Define
\[
\bar L_j
=
\frac1W\sum_{t=T-W+1}^{T}L_j^{(t)},
\]
\[
u_j
=
1-\left|\tanh\frac{\bar L_j}{2}\right|,
\]
and
\[
o_j
=
\frac1{W-1}
\sum_{t=T-W+2}^{T}
\ind\!\left[
\sgn L_j^{(t)}
\ne
\sgn L_j^{(t-1)}
\right].
\]

Let $d_j$ be the Tanner-graph distance, measured on the immutable physical
graph $G_0$, from physical variable $j$ to the nearest residual check in
$\mathcal R$. The reference failure weight is
\[
\boxed{
\omega_j
=
\ind[d_j\le D]
\,\rho_{\mathrm{dist}}^{\,d_j}
\frac{\eta_u u_j+\eta_o o_j}{\eta_u+\eta_o}.
}
\]
If $\mathcal R=\varnothing$, the projected physical hard decision already
satisfies $H_0\hat x=s_0$, so the decoder must have terminated before this
subroutine.

For qDither, $\mathcal H_L$ is taken from the final unsuccessful internal
chain only. Histories from distinct qDither chains are not concatenated.

\subsection{Selecting the suspicious physical set $U$}

Two selection modes are defined.

\paragraph{Top-$K$ mode.}
\[
U
=
\TopK_{j\in\Vphys}\omega_j,
\]
retaining at most $K_U$ variables with positive weight. Ties are broken by
ascending physical variable index.

\paragraph{Threshold mode.}
\[
U
=
\{j\in\Vphys:\omega_j\ge\tau_U\}.
\]
The implementation option is
\[
\texttt{U\_selection}\in
\{\texttt{top\_k},\texttt{threshold}\}.
\]

\section{Extending failure activity to auxiliary variables}

Graph factorization may already have introduced auxiliary variables. A current
graph variable $v\in V^{(r)}$ is considered related to the suspicious physical
set if
\[
\supp P(v)\cap U\ne\varnothing.
\]
Define
\[
U^{(r)}
=
\left\{
v\in V^{(r)}:
\supp P(v)\cap U\ne\varnothing
\right\}.
\]

For weighted-cycle scoring, extend the physical weights to every current graph
variable by
\[
\boxed{
\widetilde\omega_v
=
\begin{cases}
\max_{j\in\supp P(v)}\omega_j,
& \supp P(v)\ne\varnothing,\\
0,&\supp P(v)=\varnothing.
\end{cases}
}
\]
The maximum rule is part of the reference specification. It prevents a
large-support auxiliary variable from receiving a larger score merely because
its support contains many physical variables.

\section{Biclique-discovery subroutine}

\subsection{Relevant local graph}

Given $U^{(r)}$, define
\[
\boxed{
\Crel
=
N_{G^{(r)}}(U^{(r)})
}
\]
and, for bookkeeping,
\[
\Vrel
=
N_{G^{(r)}}(\Crel).
\]
The role of $\Vrel$ is only to describe the induced local graph
\[
G_{\mathrm{rel}}^{(r)}
=
G^{(r)}[\Vrel,\Crel].
\]
It is not used redundantly inside the neighborhood definitions below.

\subsection{Check-pair seeds}

For every unordered pair
\[
a,b\in\Crel,\qquad a<b,
\]
compute
\[
\boxed{
S_{ab}
=
N_{G^{(r)}}(a)\cap N_{G^{(r)}}(b).
}
\]
Discard the pair unless
\[
|S_{ab}|\ge2
\]
and
\[
S_{ab}\cap U^{(r)}\ne\varnothing.
\]

For every surviving pair define the check closure
\[
\boxed{
C_{ab}
=
\left\{
c\in\Crel:
S_{ab}\subseteq N_{G^{(r)}}(c)
\right\}.
}
\]
Because $S_{ab}\cap U^{(r)}\ne\varnothing$, any check in the complete graph
that contains all of $S_{ab}$ is automatically adjacent to $U^{(r)}$ and hence
belongs to $\Crel$. Therefore this local closure is also the complete check
closure for this candidate.

Moreover, because $a,b\in C_{ab}$,
\[
\bigcap_{c\in C_{ab}}N(c)
\subseteq
N(a)\cap N(b)
=
S_{ab},
\]
while the definition of $C_{ab}$ gives the opposite inclusion. Hence
\[
\boxed{
S_{ab}
=
\bigcap_{c\in C_{ab}}N(c).
}
\]
Thus the pair seed produces a closed biclique
\[
B_{ab}=(S_{ab},C_{ab}).
\]

Identical pairs $(S_{ab},C_{ab})$ generated by different check seeds are
deduplicated using sorted variable and check identifier lists.

\subsection{Output}

The output of biclique discovery is
\[
\boxed{
\Bcal(G^{(r)},U)
=
\operatorname{unique}
\left\{
(S_{ab},C_{ab})
\right\}.
}
\]

\begin{algorithm}[H]
\caption{\textsc{DiscoverBicliques}}
\begin{algorithmic}[1]
\Require Current graph $G$, support map $P$, suspicious physical set $U$
\Ensure Unique biclique set $\Bcal$
\State $U_G\gets\{v:\supp P(v)\cap U\ne\varnothing\}$
\State $\Crel\gets N_G(U_G)$
\State $\Bcal\gets\varnothing$
\ForAll{unordered pairs $a<b$ in $\Crel$}
    \State $S\gets N_G(a)\cap N_G(b)$
    \If{$|S|<2$ or $S\cap U_G=\varnothing$}
        \State \textbf{continue}
    \EndIf
    \State $C\gets\{c\in\Crel:S\subseteq N_G(c)\}$
    \State insert canonical pair $(S,C)$ into $\Bcal$
\EndFor
\State \Return $\Bcal$
\end{algorithmic}
\end{algorithm}

An implementation may generate only check pairs that share variables by
accumulating pairs from variable-to-check adjacency. This optimization must
produce exactly the same mathematical output as the algorithm above.

\section{Exact biclique factorization}

Let
\[
B=(S,C)\in\Bcal,
\qquad
|S|\ge2,\quad |C|\ge2,
\]
be a biclique in the current graph $G$. Introduce a fresh auxiliary variable
$y_B$ and a fresh auxiliary check $q_B$. Add the defining equation
\[
\boxed{
z_{y_B}
\oplus
\bigoplus_{v\in S}z_v
=0.
}
\]
Equivalently,
\[
N(q_B)=S\cup\{y_B\},
\qquad
s_{q_B}=0.
\]

For every check $c\in C$, replace
\[
N(c)
\]
by
\[
\boxed{
N'(c)
=
\left(N(c)\setminus S\right)\cup\{y_B\}.
}
\]
All other check neighborhoods remain unchanged.

The new auxiliary variable support is
\[
\boxed{
P(y_B)
=
\bigoplus_{v\in S}P(v).
}
\]
Its immutable base prior is
\[
\lambda_{y_B}^{\mathrm{base}}=0.
\]

The new check enforces $z_{y_B}=\oplus_{v\in S}z_v$; substituting this equality
into every transformed check reproduces the original check exactly. Therefore
the transformation preserves the physical solution set one-to-one. This is
the basic AVN/ACN construction used for all-ones submatrices by Shen et
al.~\cite{shen} and the biclique rewiring principle used by GARI~\cite{gari}.

If $p=|C|$ and $q=|S|$, the biclique contains $pq$ internal edges and
\[
\binom p2\binom q2
\]
internal 4-cycles. The transformed repeated-parity subgraph uses $p+q+1$
edges. Candidate A intentionally does not include edge-count, maximum-degree,
or auxiliary-count penalties in its selection score; these quantities should
still be logged as diagnostics.

\section{Failure-weighted net 4-cycle score}

For any current graph $G=(V,C,E)$, define
\[
\boxed{
\Phi(G;\omega)
=
\sum_{\{a,b\}\subseteq C}
\;
\sum_{\{u,v\}\subseteq N(a)\cap N(b)}
\widetilde\omega_u\widetilde\omega_v.
}
\]
Each unordered pair of checks and unordered pair of shared variables
corresponds to exactly one 4-cycle, so every 4-cycle is counted exactly once.

For candidate biclique $B$, let
\[
G_B=T_B(G)
\]
be the trial graph produced by exact biclique factorization. Define
\[
\boxed{
\Delta\Phi^{\mathrm{net}}(B)
=
\Phi(G;\omega)-\Phi(G_B;\omega).
}
\]
A positive value means that the transformation reduces the current
failure-weighted 4-cycle score after accounting for both removed and newly
created 4-cycles.

\section{Factorization subroutine}

\subsection{Policy A: adaptive net-cycle selection}

At each factorization step, compute the complete current biclique set
$\Bcal(G,U)$. For every $B\in\Bcal$, build the trial transformed graph $G_B$
and compute $\Delta\Phi^{\mathrm{net}}(B)$. Select
\[
\boxed{
B^\star
=
\argmax_{B\in\Bcal}
\Delta\Phi^{\mathrm{net}}(B).
}
\]
Ties are broken deterministically by lexicographic order of the sorted
variable-ID list $S$, followed by the sorted check-ID list $C$.

If
\[
\boxed{
\Delta\Phi^{\mathrm{net}}(B^\star)\le0,
}
\]
no transformation is applied and the factorization subroutine terminates
early. Thus larger positive score is better.

\subsection{Policy B: Shen-style internal-cycle baseline}

For every current biclique $B=(S,C)$, compute
\[
\boxed{
\delta(B)
=
\binom{|S|}{2}
\binom{|C|}{2}.
}
\]
Select
\[
\boxed{
B^\star
=
\argmax_{B\in\Bcal}\delta(B).
}
\]
The same deterministic lexicographic tie rule is used. A candidate in
$\Bcal$ always has $|S|\ge2$ and $|C|\ge2$, hence $\delta(B)\ge1$; this policy
applies the selected transformation whenever $\Bcal\ne\varnothing$.

This is a \emph{Shen-style selection baseline}, not a reproduction of every
detail of Shen et al.'s optimized 4-cycle-removal procedure. In particular,
the present baseline does not include their additional degree-1/degree-2
cleanup rules. The purpose of policy B is to isolate the effect of selecting
the biclique with the largest number of internal 4-cycles.

\subsection{Multiple transformations and overlapping bicliques}

Let $n_{\mathrm{fact}}$ be the maximum number of accepted factorization
operations performed in one graph-refactorization call.

Transformations are strictly sequential. After every accepted transform, the
actual graph is replaced by the transformed graph, the support map and base
prior are extended, the biclique candidate set is discarded, biclique
discovery is rerun on the new graph, and candidate scores are recomputed.
Therefore overlapping bicliques are never treated as independent static
objects. A candidate that was valid before a transformation has no standing
after the graph is modified unless it is rediscovered in the updated graph.

The physical failure weights $\omega$ and suspicious physical set $U$ remain
fixed throughout the $n_{\mathrm{fact}}$ inner transformations because no BP
run occurs between them. Their extension $\widetilde\omega$ to newly introduced
auxiliary variables is recomputed from the updated physical-support map.

\begin{algorithm}[H]
\caption{\textsc{FactorizeGraph}}
\begin{algorithmic}[1]
\Require $G,\lambda^{\mathrm{base}},P,\omega,U,n_{\mathrm{fact}}$, policy
\Ensure Transformed $G,\lambda^{\mathrm{base}},P$, number $n_{\mathrm{applied}}$
\State $n_{\mathrm{applied}}\gets0$
\For{$t=1,\ldots,n_{\mathrm{fact}}$}
    \State $\Bcal\gets\textsc{DiscoverBicliques}(G,P,U)$
    \If{$\Bcal=\varnothing$}
        \State \textbf{break}
    \EndIf
    \If{policy = \texttt{adaptive\_cycle}}
        \ForAll{$B\in\Bcal$}
            \State form trial $G_B\gets T_B(G)$
            \State $\Delta_B\gets\Phi(G;\omega)-\Phi(G_B;\omega)$
        \EndFor
        \State $B^\star\gets\argmax_B\Delta_B$
        \If{$\Delta_{B^\star}\le0$}
            \State \textbf{break}
        \EndIf
    \Else
        \State $\delta_B\gets\binom{|S_B|}{2}\binom{|C_B|}{2}$ for all $B\in\Bcal$
        \State $B^\star\gets\argmax_B\delta_B$
    \EndIf
    \State Apply $T_{B^\star}$ exactly to $G,\lambda^{\mathrm{base}},P$
    \State $n_{\mathrm{applied}}\gets n_{\mathrm{applied}}+1$
\EndFor
\State \Return updated state and $n_{\mathrm{applied}}$
\end{algorithmic}
\end{algorithm}

Shen et al. similarly apply a greedy 4-cycle-removal operation and continue on
the updated matrix rather than treating initially detected overlapping
all-ones blocks as independent transformations~\cite{shen}. The adaptive
procedure here differs in its selection objective and in stopping after a
configurable number of transformations.

\section{Marginal handoff after graph transformation}
\label{sec:handoff}

The graph-relay algorithm inherits the previous BP state only as an
\emph{initialization}; it does not overwrite the immutable unary factors.

Let $L_v^{\mathrm{parent}}$ be the final marginal of the failed BP call for
every variable that existed before factorization. Define the relay
initialization vector $q^{\mathrm{relay}}$ on the transformed graph. For every
surviving existing variable,
\[
q_v^{\mathrm{relay}}
=
L_v^{\mathrm{parent}}.
\]

For a newly introduced auxiliary variable
\[
y=\bigoplus_{v\in S}z_v,
\]
compute an approximate parity LLR from the currently available relay marginals:
\[
\boxed{
q_y^{\mathrm{relay}}
=
2\,\operatorname{atanh}
\left(
\prod_{v\in S}
\tanh\frac{q_v^{\mathrm{relay}}}{2}
\right).
}
\]
Numerically, the product must be clipped to
\[
[-1+\epsilon_{\mathrm{atanh}},\,1-\epsilon_{\mathrm{atanh}}].
\]
This formula is an independence approximation used only to initialize the new
auxiliary message state; it is not treated as a new physical prior.

If several factorizations are applied in one call, the new auxiliary marginal
is computed immediately after each accepted transform so that it is available
if a later transform uses that auxiliary variable.

For ordinary parallel or serial BP, initialize
\[
m_{v\to a}^{(0)}
=
q_v^{\mathrm{relay}}
\]
on every edge of the transformed graph. Starting with the first actual BP
iteration, all variable updates use the immutable $\lambda_v^{\mathrm{base}}$,
not $q_v^{\mathrm{relay}}$. Thus marginal inheritance changes the BP trajectory
but does not change the exact factorization.

\section{Top-level adaptive graph-relay decoder}

Let $R_{\mathrm{graph}}$ denote the maximum number of
factorization--redecode rounds \emph{after} the initial BP call. Hence the
maximum number of BP calls is $R_{\mathrm{graph}}+1$.

\begin{algorithm}[H]
\caption{\textsc{AdaptiveGraphRelayDecode}}
\begin{algorithmic}[1]
\Require $H_0,s_0,p$, BP options, failure-weight options,
         $U$-selection options, factorization policy,
         $n_{\mathrm{fact}},R_{\mathrm{graph}}$
\Ensure Valid physical correction or declared failure
\State Build $G\gets G_0$, $P(x_j)\gets e_j$,
       $\lambda_j^{\mathrm{base}}\gets\log((1-p_j)/p_j)$
\State $q^{\mathrm{init}}\gets\lambda^{\mathrm{base}}$
\For{$r=0,\ldots,R_{\mathrm{graph}}$}
    \State $R_{\mathrm{BP}}\gets
    \textsc{RunBP}(G,\lambda^{\mathrm{base}},q^{\mathrm{init}},\Theta_{\mathrm{BP}})$
    \If{$R_{\mathrm{BP}}$ contains $\hat x$ with $H_0\hat x=s_0$}
        \State \Return $\hat x$
    \EndIf
    \If{$r=R_{\mathrm{graph}}$}
        \State \textbf{break}
    \EndIf
    \State $(\omega,U)\gets
    \textsc{ComputeFailureWeights}(R_{\mathrm{BP}},H_0,s_0,\Theta_\omega)$
    \If{$U=\varnothing$}
        \State \textbf{break}
    \EndIf
    \State Save $L^{\mathrm{parent}}\gets R_{\mathrm{BP}}.L^{\mathrm{final}}$
    \State $(G,\lambda^{\mathrm{base}},P,n_{\mathrm{applied}})
    \gets\textsc{FactorizeGraph}
    (G,\lambda^{\mathrm{base}},P,\omega,U,n_{\mathrm{fact}},\text{policy})$
    \If{$n_{\mathrm{applied}}=0$}
        \State \textbf{break}
    \EndIf
    \State $q^{\mathrm{init}}\gets
    \textsc{TransferMarginals}(L^{\mathrm{parent}},\text{factorization log})$
\EndFor
\State \Return declared failure
\end{algorithmic}
\end{algorithm}

\section{Subroutine contracts}

\begin{longtable}{@{}p{34mm}p{48mm}p{48mm}p{33mm}@{}}
\toprule
Subroutine & Input & Output & Main options \\
\midrule
\endhead
\textsc{RunBP}
&
$G,\lambda^{\mathrm{base}},q^{\mathrm{init}}$
&
convergence flag, hard decision, final marginals, trailing marginal history
&
BP variant, iteration budget, Min-Sum normalization, serial order, qDither parameters
\\[1mm]

\textsc{ComputeFailureWeights}
&
failed BP result, immutable $H_0,s_0,G_0$
&
physical $\omega$, suspicious physical set $U$
&
history window, residual radius, distance decay, uncertainty/oscillation weights, $U$-selection mode
\\[1mm]

\textsc{DiscoverBicliques}
&
current $G$, support map $P$, physical $U$
&
unique current bicliques $\Bcal$
&
pair seeds only in version 1
\\[1mm]

\textsc{ApplyBicliqueFactorization}
&
current graph state, one biclique $B=(S,C)$
&
exact transformed graph, extended base prior, extended support map
&
none in version 1
\\[1mm]

\textsc{FactorizeGraph}
&
current graph state, $\omega,U,n_{\mathrm{fact}}$
&
transformed graph state and number of accepted transforms
&
\texttt{adaptive\_cycle} or \texttt{shen\_cycle\_count}
\\[1mm]

\textsc{TransferMarginals}
&
parent final marginals and accepted transform sequence
&
initialization vector on transformed graph
&
parity-LLR clipping parameter
\\
\bottomrule
\end{longtable}

\section{Initial parameter set}

The values below are implementation defaults for the first study; they are not
claimed to be optimized.

\begin{longtable}{@{}p{46mm}p{31mm}p{84mm}@{}}
\toprule
Parameter & Initial default & Meaning \\
\midrule
\endhead
\texttt{bp\_variant}
& \texttt{parallel}
& \texttt{parallel}, \texttt{serial}, or \texttt{qdither}. \\
\texttt{bp\_iterations}
& 50
& Maximum iterations for one ordinary parallel/serial BP call. \\
\texttt{ms\_scale}
& 1.0
& Min-Sum normalization $\alpha_{\mathrm{MS}}$. \\
\texttt{serial\_order}
& \texttt{random\_per\_iteration}
& Horizontal serial check order. \\
\texttt{history\_window}
& 8
& Number of final marginal iterations used in $u_j,o_j$. \\
\texttt{residual\_radius} $D$
& 2
& Maximum original-graph distance receiving nonzero $\omega_j$. \\
\texttt{distance\_decay} $\rho_{\mathrm{dist}}$
& 0.5
& Multiplicative residual-distance decay. \\
\texttt{uncertainty\_weight} $\eta_u$
& 1
& Relative contribution of marginal uncertainty. \\
\texttt{oscillation\_weight} $\eta_o$
& 1
& Relative contribution of sign oscillation. \\
\texttt{U\_selection}
& \texttt{top\_k}
& Suspicious-variable selection rule. \\
\texttt{U\_top\_k} $K_U$
& 32
& Maximum suspicious physical variables in top-$K$ mode. \\
\texttt{U\_threshold} $\tau_U$
& 0.5
& Used only in threshold mode. \\
\texttt{factorization\_policy}
& \texttt{adaptive\_cycle}
& Candidate A or Shen-style candidate B. \\
\texttt{n\_fact}
& 1
& Maximum accepted factorizations between two BP calls. \\
\texttt{graph\_rounds} $R_{\mathrm{graph}}$
& 4
& Maximum factorization--redecode rounds after initial BP. \\
\texttt{atanh\_epsilon}
& $10^{-12}$
& Numerical clip for auxiliary parity-LLR initialization. \\
\texttt{qdither\_chains}
& implementation value
& Number of qDither Phase-II chains. \\
\texttt{qdither\_iterations\_per\_chain}
& implementation value
& Iteration budget $b$ for each qDither chain. \\
\texttt{qdither\_alpha}, \texttt{qdither\_beta}
& experiment-specific
& Uniform mixing range $[\alpha_D,\beta_D]$ in qDither. \\
\texttt{qdither\_rho}
& experiment-specific
& Bernoulli probability for enhanced RSF in qDither. \\
\texttt{qdither\_handoff}
& \texttt{paper}
& Faithful paper initialization or explicit graph-warm hybrid. \\
\bottomrule
\end{longtable}

For qDither, paper-specific numerical values depend on the benchmark and should
not be hard-coded as universal defaults~\cite{qdither}.

\section{Complexity of the initial design}

Let
\[
R=|\Crel|,
\qquad
d_v=\max_{v\in V}|N(v)|,
\qquad
d_c=\max_{c\in C}|N(c)|.
\]
Naively testing all unordered check pairs requires
\[
\binom R2=O(R^2)
\]
pair seeds.

In a sparse implementation, pairs can instead be accumulated from each
variable $v$ by enumerating
\[
\binom{|N(v)\cap\Crel|}{2}
\]
check pairs. The total pair-incidence work is
\[
O\!\left(
\sum_{v\in\Vrel}
\binom{|N(v)\cap\Crel|}{2}
\right),
\]
which is linear in the local graph size for bounded-degree qLDPC graphs.

Policy B requires only biclique sizes after discovery. Policy A is more
expensive because every candidate requires a trial transform and an exact net
weighted-4-cycle delta. A correct implementation need not copy the whole graph
for each trial: only check pairs whose neighborhoods change under $T_B$
contribute to $\Delta\Phi^{\mathrm{net}}(B)$. Local incremental evaluation is
therefore an implementation optimization that preserves the reference
definition.

The factorization stage performs at most $n_{\mathrm{fact}}$ accepted
transforms per graph-relay round and recomputes candidates after every accepted
transform.

\section{Required validation invariants}

The implementation should assert or unit-test the following properties.

\paragraph{Exact factorization.}
For small random instances, enumerate all physical $x$ and verify
\[
H_0x=s_0
\iff
\exists!\,z_{\mathrm{aux}}:\;H^{(r)}z=s^{(r)}.
\]

\paragraph{Support-map correctness.}
For every auxiliary variable, verify that its value equals $P(v)^Tx$ on
enumerated assignments.

\paragraph{Biclique correctness.}
Every reported $B=(S,C)$ must satisfy
\[
S\subseteq N(c)
\quad\forall c\in C,
\]
\[
|S|\ge2,\qquad |C|\ge2,
\]
and the closure identities specified in the discovery section.

\paragraph{Cycle-score correctness.}
For small graphs, compare $\Phi$ and $\Delta\Phi^{\mathrm{net}}$ against explicit
4-cycle enumeration.

\paragraph{Overlap handling.}
Construct overlapping bicliques, apply one transform, and verify that the next
candidate set is regenerated from the transformed graph rather than reused
from the stale candidate list.

\paragraph{Prior invariance.}
Physical base priors must never change under graph factorization; newly
introduced auxiliary variables must have zero base LLR.

\paragraph{Warm-start semantics.}
Inherited marginals may initialize edge messages but must not replace the
immutable base unary factors in ordinary parallel/serial BP.

\paragraph{Original-problem validation.}
Every declared successful correction must satisfy
\[
H_0\hat x=s_0
\]
after projection to physical variables.

\section{Logical-ambiguity audit}

The following points are fixed explicitly in this specification because they
would otherwise be ambiguous in implementation.

\begin{enumerate}[leftmargin=2em]
\item \textbf{Sign of policy-A score.}
      The score is before minus after. Larger is better; a non-positive best
      score causes no transform.

\item \textbf{Meaning of $n_{\mathrm{fact}}$.}
      It counts accepted transformations between two BP calls, not candidate
      evaluations and not graph-relay rounds.

\item \textbf{Overlapping bicliques.}
      They are not batch-applied. Candidate discovery and scoring are rerun
      after every accepted transformation.

\item \textbf{Failure weights during several inner factorizations.}
      The physical $\omega$ and $U$ from the preceding failed BP call remain
      fixed until BP is run again. Only their extension to new auxiliaries is
      recomputed.

\item \textbf{Search depth.}
      Version 1 uses only unordered check-pair seeds. Three-check and larger
      seed sets are outside this implementation.

\item \textbf{Relevant checks.}
      They are exactly the current checks adjacent to current variables whose
      physical support intersects $U$.

\item \textbf{Candidate closure.}
      $C_{ab}$ is the complete closure of $S_{ab}$, not merely the seed pair
      $\{a,b\}$.

\item \textbf{Auxiliary priors.}
      Auxiliary variables have zero immutable unary LLR. Parent marginals are
      initialization state, not new physical priors.

\item \textbf{qDither comparison.}
      qDither is not merely a message-update schedule. It changes the iterative
      dynamics through enhanced RSF and chained randomized bias inheritance.
      Therefore the implementation option is named \texttt{bp\_variant}, not
      \texttt{schedule}. Paper-faithful and graph-warm qDither initialization
      are reported separately.

\item \textbf{Policy B terminology.}
      The implemented baseline is Shen-style biclique selection by
      $\binom{|S|}{2}\binom{|C|}{2}$. Unless the degree-1/degree-2 cleanup from
      Shen et al. is also implemented, it must not be called an exact
      reproduction of their complete optimized removal algorithm.

\item \textbf{Graph-round count.}
      $R_{\mathrm{graph}}$ counts factorization--redecode rounds after the
      initial BP call. The total number of BP calls is at most
      $R_{\mathrm{graph}}+1$.

\item \textbf{Physical diagnostic basis.}
      The residual syndrome used to form $\omega$ is always computed with
      $H_0,s_0$, never with the changing augmented check basis.
\end{enumerate}

\section{Initial comparison matrix}

The implementation supports two independent experimental axes. First,
factorization selection:
\[
\texttt{adaptive\_cycle}
\quad\text{vs.}\quad
\texttt{shen\_cycle\_count}.
\]
Second, BP dynamics:
\[
\texttt{parallel}
\quad\text{vs.}\quad
\texttt{serial}
\quad\text{vs.}\quad
\texttt{qdither}.
\]
The full adaptive decoder comparison is therefore a $2\times3$ design before
varying $n_{\mathrm{fact}}$, graph rounds, or suspicious-variable selection.

This comparison is end-to-end. Because $\omega$, $U$, and therefore the
selected graph transformations depend on the preceding BP trajectory, changing
the BP variant can also change the graph sequence. A separate mechanistic
ablation that freezes a transformation sequence would be needed to isolate
only the effect of the BP update schedule on a fixed graph.

\section{Summary of the reference version}

The first implementation is fully determined by the following sequence:
\[
\boxed{
\begin{aligned}
&\text{Run the selected BP variant on }G^{(r)}.\\
&\text{If it fails, compute physical }\omega\text{ and }U.\\
&\text{Repeat up to }n_{\mathrm{fact}}\text{ times:}\\
&\qquad
\text{discover all pair-seeded closed bicliques on the current graph;}\\
&\qquad
\text{select one biclique by policy A or policy B;}\\
&\qquad
\text{apply one exact AVN/ACN factorization;}\\
&\qquad
\text{discard and recompute the candidate set.}\\
&\text{Transfer the failed BP marginals as an initialization state.}\\
&\text{Run BP again on the transformed graph.}\\
&\text{Repeat for at most }R_{\mathrm{graph}}\text{ graph-relay rounds.}
\end{aligned}
}
\]

No batch application of stale bicliques, no higher-order check-set seed search,
and no edge/degree/auxiliary penalty is part of candidate A in this initial
version.

\begin{thebibliography}{9}

\bibitem{gari}
A.~S.~Maan, F.~M.~Garcia Herrero, A.~Paler, and V.~Savin,
``Decoding correlated errors in quantum LDPC codes,''
\emph{Nature Communications}, vol.~17, 3965, 2026.
\url{https://doi.org/10.1038/s41467-026-70556-3}.

\bibitem{shen}
Y.~Shen, Z.~Li, Y.~Ren, E.~Boutillon, A.~Balatsoukas-Stimming,
C.~Zhang, X.~You, and A.~Burg,
``Toward Universal Belief Propagation Decoding for Short Binary Block Codes,''
\emph{IEEE Journal on Selected Areas in Communications},
vol.~43, no.~4, pp.~1135--1152, 2025.
\url{https://doi.org/10.1109/JSAC.2025.3536505}.

\bibitem{qdither}
Z.~Cai, M.~A.~Jarkas, B.~Zhao, and W.~J.~Gross,
``Dithered Belief Propagation for Real-Time Quantum Memory Decoding,''
in \emph{2026 IEEE International Conference on Quantum Computing and
Engineering (QCE)}, 2026.
DOI: 10.1109/QCE68830.2026.00100.

\end{thebibliography}

\end{document}
